\section{Introducción}
La arquitectura de seguridad facilita la evaluación de la integración de los procesos, la información las aplicaciones y las tecnologías de una organización con los requisitos de protección. Bajo este marco se desarrollan diferentes escenarios de negocio y se elige la gestión y despacho de flota de TransMilenio S.A como caso de estudio procedente del Centro de Control Integrado (CCI).

Para este asunto se realiza un estudio de arquitectura empresarial usando el estándar TOGAF, abordando aspectos de negocio, información, aplicación y tecnología. Por último se reconocen y valoran los activos de información más importantes y los riesgos de seguridad que sobre ellos recaen con la finalidad de establecer cómo podrían llegar a afectar el funcionamiento y el cumplimiento de los objetivos estratégicos de la organización.